\documentclass[11pt]{article}

% A deliberately conventional preprint source: easy to import, diff, and restyle.
\usepackage[T1]{fontenc}
\usepackage[utf8]{inputenc}
\usepackage{lmodern}
\usepackage{microtype}
\usepackage[margin=1in]{geometry}
\usepackage{amsmath}
\usepackage{amssymb}
\usepackage{booktabs}
\usepackage{array}
\usepackage{tabularx}
\usepackage[section]{placeins}
\usepackage{graphicx}
\usepackage[dvipsnames]{xcolor}
\usepackage[numbers,sort&compress]{natbib}
\usepackage[
  colorlinks=true,
  linkcolor=MidnightBlue,
  citecolor=MidnightBlue,
  urlcolor=MidnightBlue,
  pdftitle={Molecular Design Software Is Now Free. Molecular Prediction Is Not.},
  pdfauthor={Rafal P. Wiewiora and Woody Sherman},
  pdfsubject={Molarium as a case study in generated scientific infrastructure},
  pdfkeywords={molecular design, agentic software, WebGPU, scientific validation, computational chemistry}
]{hyperref}

\graphicspath{{figures/}}
\newcommand{\molarium}{\textnormal{\textsc{Molarium}}}
\newcommand{\software}[1]{\texttt{#1}}
\setlength{\parindent}{0pt}
\setlength{\parskip}{0.48em}
\renewcommand{\arraystretch}{1.12}

\title{\textbf{Molecular Design Software Is Now Free.\\Molecular Prediction Is Not.}\\[0.55em]
\large \molarium{} as a case study in generated scientific infrastructure}
\author{Rafal P. Wiewiora \and Woody Sherman\\[0.35em]
\normalsize PsiThera, Watertown, Massachusetts, USA}
\date{Perspective draft --- August 2026}

\begin{document}
\maketitle

\begin{abstract}
Computational molecular design has historically bundled scientific models, numerical implementations, interfaces, workflow plumbing, and expert judgment into the same software product.
Coding agents are beginning to separate those layers.
They can increasingly generate interfaces, integrations, analyses, workflows, and bounded pieces of numerical software on demand.
``Free'' is not literal: inference, review, validation, security, operation, and maintenance remain costly.
The change is a collapse in the marginal time and cost required to reproduce parts of the implementation and orchestration layer.

We argue that value in computational drug discovery is separating into three layers: implementation and orchestration, which are becoming generated infrastructure; scientific models, which remain differentiated by accuracy and validation; and proprietary discovery intelligence, including experimental data, project history, design rationale, expert judgment, and organizational objectives.
Agents are already making computational chemistry easier to access.
They do not automatically make it more predictive.

\molarium{} provides an illustrative case rather than a controlled productivity study.
Over an observed interval of several days, an expert chemist working with GPT-5.6 Sol produced a local-first browser workbench and browser-native implementations of OpenFF Sage forces, OBC2/ACE solvent, homogeneous-replica molecular dynamics with SHAKE/RATTLE, and GPU-resident ANI-2x descriptors and forces.
The project preserved rather than regenerated RDKit, OpenMM, force-field parameters, official model weights, and accumulated scientific evidence.
Reference comparisons support implementation parity for selected cases; workload-specific timings support acceleration for sufficiently large or batched browser work; and an early conformer experiment supports added exploration, not generally better conformer prediction.
\molarium{} illustrates both sides of the thesis: substantial molecular-design software can now be generated during the scientific conversation itself.
Scientific truth cannot.
\end{abstract}

\noindent\textbf{Keywords:}
agentic software development; computational chemistry; molecular design; WebGPU; scientific validation; institutional knowledge

\section{The implementation layer is becoming generated infrastructure}
\label{sec:generated}

``Molecular design software is now free'' is useful only as a provocation.
Software is not free in the literal economic sense.
Compute, model inference, review, testing, security, support, and maintenance all consume scarce resources.
Nor is every scientific method open, portable, or appropriate for local execution.
The consequential change is narrower: \emph{implementation is becoming less scarce}.

For most of the history of computational drug discovery, a request for a new internal tool entered a queue.
Someone had to translate the scientific need into a specification, implement the data model and interface, connect external programs, convert formats, write tests, deploy the result, and maintain it.
Even when the underlying algorithm was public, the surrounding work could delay a useful workflow for months.

Coding agents change the latency between a scientific request and an executable artifact.
Natural-language instructions, examples, test cases, reference outputs, and interactive feedback can jointly serve as a specification.
An agent can inspect a repository, write and execute code, compare outputs, diagnose failures, and revise the implementation.
Repository-scale benchmarks have made this a serious software-engineering capability rather than a code-completion curiosity~\cite{jimenez2024swebench}.
Recent systems sharpen both the opportunity and the limitation.
Ayg\"un and colleagues showed that an AI system combining language models with search over candidate programs can generate high-performing empirical scientific software across several domains~\cite{aygun2026empirical}.
By contrast, the AutoMat benchmark found that coding agents still struggled to reconstruct and support claims from underspecified computational materials-science papers; its best evaluated setting succeeded on 54.1\% of expert-curated tasks~\cite{huang2026automat}.
In chemistry, ChemCrow and Coscientist have shown how language-model agents can orchestrate specialist tools, search, code, and experimental interfaces~\cite{bran2024chemcrow,boiko2023coscientist}.

\molarium{} adds a different kind of evidence.
The coding agent did not only decide which established program to call.
Under expert direction, it implemented missing browser-specific scientific infrastructure, including substantial numerical kernels.
The case does not establish a universal productivity multiplier, and active development time has not yet been reconstructed from the project archive.
It shows only that a meaningful, bounded class of molecular-design implementation can now be created while the scientific question is still being formed.

``Free'' can therefore be decomposed into five claims:
\begin{enumerate}
  \item \textbf{Free to use.} Original application code can be distributed without a license fee, and supported local calculations need not require a paid service.
  \item \textbf{Free to rebuild.} A sufficiently explicit method may be reimplemented from equations, parameters, tests, and reference values rather than awaited as a product feature.
  \item \textbf{Free to preserve.} The generated method, its tests, and its provenance can remain as inspectable code rather than disappearing into a transient conversation.
  \item \textbf{Free to compose.} Interfaces and workflows can be built around the current scientific question rather than forcing the question into a fixed application.
  \item \textbf{Not free to validate or operate at scale.} Scientific credibility, security, maintenance, heterogeneous hardware support, and dependable operation remain expensive.
\end{enumerate}

The qualification is the point.
Lower implementation cost is important because it exposes what was never reducible to implementation in the first place.

\section{Where value moves}
\label{sec:value}

Computational molecular design can be viewed as three coupled layers (Figure~\ref{fig:layers}).

\begin{figure}[htbp]
  \centering
  \includegraphics[width=\linewidth]{three_layers.pdf}
  \caption{Three layers of value in computational molecular science.
  Implementation and orchestration are becoming increasingly reproducible by agents.
  Scientific models remain differentiated by accuracy, domain coverage, and validation.
  Proprietary discovery intelligence may become the deepest moat because it combines evidence, context, and judgment that are difficult to imitate.}
  \label{fig:layers}
\end{figure}

The first layer has historically carried substantial commercial value.
Companies and internal platform teams implemented algorithms, exposed command-line programs to chemists, built graphical interfaces, converted formats, integrated third-party tools, maintained workflow plumbing, and produced bespoke analysis views.
Agents can now perform a remarkable fraction of this work on demand.
A moat defined primarily as ``we integrated the tools and arranged the interface well'' is therefore becoming easier to reproduce.
Mature applications still embody years of edge cases, compatibility work, user trust, support, and operational discipline; those assets should not be confused with a screen that can be regenerated.

That does not make scientific-software companies irrelevant.
It changes what customers are paying for.
The more durable assets are validated models, high-quality data, proprietary parameterization, curated benchmarks, secure deployment, accountable scientific support, and organization-specific workflows whose decisions have been tested.
The application remains important, but it increasingly expresses the moat rather than constituting it.

The boundaries among the three layers are porous.
Protonation policy, charge assignment, precision, integrator choice, stopping rules, and default visualization may look like implementation details but can alter a scientific result.
``Implementation layer'' therefore does not mean ``scientifically neutral.''
It means the part that can, in principle, be reconstructed from an explicit contract and subjected to falsifiable acceptance tests.

The distinction between \emph{tool space} and \emph{model space} is equally important.
Most current agents operate in tool space: they choose and coordinate docking, molecular dynamics, free-energy calculations, quantum chemistry, descriptors, or learned models.
Better orchestration can make these methods easier to use, but it inherits their representations and failure modes.
Predictive progress requires better models of molecular state, better training data, and better uncertainty---not merely a more fluent agent choosing from the existing menu~\cite{bender2021part1,bender2021part2,wigh2022representation}.

This gives computational chemistry two different revolutions.
The first is to make it accessible.
The second, harder revolution is to make it predictive.
Agents are already transforming the first.
They do not automatically solve the second.

\section{\molarium: a case study rather than a proof}
\label{sec:case}

\molarium{} is a public, local-first molecular viewer, builder, and simulation workbench~\cite{wiewiora2026molarium}.
Most supported calculations run inside a contemporary browser using WebAssembly or WebGPU.
The workbench can load small molecules and protein structures; preserve explicit topology during three-dimensional editing; prepare supported protein structures; run energies, minimization, molecular dynamics, and conformer searches; evaluate ANI-2x; and run a restricted, short-chain OpenFold model.
A dedicated Local Lab mode restricts network access for proprietary structures and verifies the files that comprise the reviewed application build.
This is an auditable technical control, not an absolute privacy guarantee: a hosted page controls the code it delivers, so the strongest boundary is a reviewed, versioned local checkout with independently verified assets.

The browser matters for two reasons.
First, it removes much of the installation and data-custody friction that has historically separated a scientist from a calculation.
WebAssembly provides a portable target for mature compiled libraries, while WebGPU provides local accelerator access~\cite{haas2017wasm,w3c2025webgpu}.
Browser molecular graphics were already well established~\cite{rose2015ngl}.
Second, the browser is an unforgiving implementation environment.
There may be no Python environment, CUDA toolkit, or persistent filesystem.
WebGPU generally offers fast single-precision arithmetic but not the full programming model or numerical capabilities of native CUDA.
These constraints make the boundary between scientific specification and implementation visible.

\subsection{The development conversation became a scientific build loop}
\label{sec:loop}

The project did not emerge from one successful prompt.
Early visual failures changed the representation of the molecule.
A benzene was wrong.
A methyl fragment formed a false bond from spatial proximity.
A trifluoromethyl group arrived flat.
A bond-order edit triggered geometry cleanup before the coupled chemical change was complete.
These observations led to explicit topology, three-dimensional fragment coordinates, staged edits, and local rather than global cleanup.

The same loop governed the numerics.
Asking whether an early evaluator was a real force field moved the project to RDKit MMFF94/UFF and prepared OpenFF Sage systems~\cite{landrum2025rdkit,riniker2015etkdg,halgren1996mmff94,rappe1992uff,boothroyd2023sage}.
A disagreement between CPU and GPU energies exposed both a unit error and an architectural misunderstanding: OpenMM Reference compiled to WebAssembly was a CPU oracle, not the WebGPU implementation~\cite{eastman2024openmm8}.
Asking whether the application was really using WebGPU led to analytic WGSL force kernels.
Asking for implicit solvent led to OBC2/ACE~\cite{onufriev2004obc,schaefer1996ace}.
Asking whether constrained dynamics were supported led to SHAKE/RATTLE~\cite{ryckaert1977shake,andersen1983rattle}.
A requested long run that stopped at a development cap exposed a product limit accidentally inherited from a test.

The recurring loop is shown in Figure~\ref{fig:loop}.

\begin{figure}[htbp]
  \centering
  \includegraphics[width=\linewidth]{build_loop.pdf}
  \caption{The Molarium scientific build loop.
  Visual, chemical, numerical, and workflow failures became implementation changes.
  Independent reference paths determined whether those changes survived, and the repository retained accepted methods with tests, benchmark boundaries, limitations, and provenance.}
  \label{fig:loop}
\end{figure}

The scientist's contribution was not a vague preference for a nicer interface.
It was the ability to notice what was chemically implausible, ask whether the engine and units were what the interface claimed, select a meaningful reference, and decide which failure mattered.
The agent supplied implementation speed.
The scientist supplied the scientific standard.

\subsection{What was generated and what was preserved}
\label{sec:generated-preserved}

The \molarium{} case is useful because the boundary is unusually explicit.
Table~\ref{tab:generated-preserved} summarizes the division.

\begin{table}[htbp]
\centering
\caption{Generated infrastructure versus preserved scientific assets in Molarium.}
\label{tab:generated-preserved}
\footnotesize
\begin{tabularx}{\linewidth}{@{}>{\raggedright\arraybackslash}X >{\raggedright\arraybackslash}X@{}}
\toprule
\textbf{Generated or substantially implemented for Molarium} & \textbf{Preserved, reused, or treated as an external scientific asset} \\
\midrule
Persistent molecular canvas and topology-aware builder & RDKit chemical perception, file handling, ETKDGv3, MMFF94, and UFF \\
Browser preparation, conformer, trajectory, and inspection workflows & OpenMM Reference as a slower numerical oracle \\
Direct WebGPU evaluation of prepared Sage bonded, electrostatic, Lennard--Jones, and OBC2/ACE terms & Published force-field equations and OpenFF parameter sets \\
Homogeneous-replica WebGPU engine with fixed-point accumulation, Langevin/Verlet integration, and replica-local SHAKE/RATTLE & STORMM as architectural inspiration and an independent upstream reference, not linked runtime code~\cite{cerutti2024stormm} \\
WebGPU construction of ANI-2x atomic environment vectors and analytic contraction of descriptor derivatives to Cartesian forces & Official TorchANI/ANI-2x architecture, ensemble weights, and native reference values~\cite{gao2020torchani,devereux2020ani2x} \\
Local-mode network controls, build manifests, egress tests, provenance views, and regression fixtures & OpenFold weights and their accumulated training and validation evidence~\cite{ahdritz2024openfold} \\
\bottomrule
\end{tabularx}
\end{table}

This is more than glue code.
The browser engines required new data layouts and algorithms.
A prepared numeric molecular system---masses, charges, bonded terms, Lennard--Jones parameters, exceptions, constraints, and optional generalized-Born data---forms a common boundary.
One execution path evaluates one system and trajectory.
A second adds a replica dimension for many conformers of one topology.
A third constructs ANI-2x descriptors and forces around preserved neural-network weights.
Each path has a deliberately separate reference route (Figure~\ref{fig:architecture}).

\begin{figure}[htbp]
  \centering
  \includegraphics[width=\linewidth]{architecture.pdf}
  \caption{One explicit chemical boundary, three browser workloads.
  Preparation ends at a declared numeric system rather than leaving the GPU kernels to infer chemistry.
  Direct Sage WebGPU, homogeneous-replica WebGPU, and ANI-2x/ONNX consume that boundary or its coordinates, and each has a deliberately separate reference path.}
  \label{fig:architecture}
\end{figure}

Yet the generated code did not regenerate the science embodied in parameter sets, trained weights, and accumulated validation.
A force-field equation can be reimplemented and checked against reference values.
A neural architecture cannot recreate the learned information in its weights or the experimental evidence that establishes where it works.
\molarium{}'s original code is MIT-licensed, while bundled libraries, force fields, models, and data retain their upstream licenses.
Generation does not erase provenance or authorship.

The present scientific scope is deliberately narrow.
The browser engines are nonperiodic and omit PME, barostats, membranes, and arbitrary OpenMM plugins.
The replica engine accepts homogeneous copies of one topology, with a current per-replica size limit.
ANI-2x is restricted to neutral, closed-shell molecules in its supported elemental domain; the browser OpenFold path handles one short chain; and arbitrary browser-native protein parameterization remains experimental.
These boundaries are part of the result, not footnotes to be hidden~\cite{wiewiora2026molarium}.

The central division is therefore:
\begin{quote}
\textbf{Equations and algorithms can often be reconstructed and checked.
Parameters, weights, data, and accumulated scientific evidence must be preserved.}
\end{quote}

\section{What \molarium{} demonstrates---and what it does not}
\label{sec:evidence}

The strongest version of the paper depends on separating three questions that software demonstrations often collapse.

\subsection{Did the implementation match its declared reference?}

For selected fixtures, the answer is encouraging.
Direct Sage WebGPU is compared with OpenMM Reference on the same prepared numeric system.
ANI-2x browser evaluation is compared with native TorchANI.
Analytic forces are checked with complete reference force vectors and finite differences.
Small fixtures test units, exceptions, constraints, replica ordering, translation invariance, and failure behavior.

On the recorded Apple M1 Pro development stack, the repository provisionally reports a maximum relative direct-Sage single-point deviation of approximately $2.4\times10^{-6}$ across its molecule panel and an aspirin force relative RMS of approximately $7.7\times10^{-7}$.
The reported aspirin OBC2/ACE case differed from OpenMM by approximately $1.3\times10^{-4}$~kcal/mol, with a force relative RMS of approximately $7.8\times10^{-7}$.
Against six native TorchANI fixtures, the browser ANI-2x path reported a maximum energy error of 0.02572~kcal/mol and a maximum relative force RMS of approximately $1.0\times10^{-5}$~\cite{wiewiora2026molarium}.

These are provisional implementation-parity results on a particular browser and GPU stack.
Publication claims should ultimately be generated from a frozen, machine-readable ledger containing the commit, fixture and model hashes, browser, adapter, driver, precision, and timing boundary.
They do not validate Sage as a model of arbitrary chemistry, independently verify upstream parameter assignment, establish ANI-2x transferability, or demonstrate cross-vendor reproducibility.
Shared parameters can allow two implementations to agree while preserving the same upstream mistake.
That is why the repository distinguishes arithmetic parity from scientific validity and lists hand-computable fixtures, finite differences, native references, force tests, dynamics tests, and a cross-GPU matrix as separate validation gates.

\subsection{Did the browser implementation make useful work faster?}

Sometimes---and the negative cases are informative.
GPU dispatch is a poor bargain for one tiny molecule.
The project's early measurements found that a single small replica could be slower than OpenMM Reference WebAssembly because dispatch and synchronization dominated the arithmetic.
Larger or batched work reversed the comparison.

On one Apple M1 Pro, the repository records a 2.09-fold throughput advantage over OpenMM Reference WebAssembly for a warm 1,000-step, 304-atom Trp-cage vacuum run; a 9.16-fold worker-time advantage for 1,231-atom ubiquitin in vacuum; and a 25.6-fold advantage for the recorded ubiquitin OBC2/ACE workload.
At equal aggregate warm integration steps, 1,024 C16 replicas reached a reported 10.3-fold advantage and 256 water27 replicas a 24.1-fold advantage over the stated Reference baseline.
Moving ANI-2x descriptor and force work into GPU-resident buffers accelerated that stage 6.77-fold relative to the previous browser path~\cite{wiewiora2026molarium}.

None of these is a native-CUDA claim.
The direct and replica baselines, integrators, and timing boundaries differ; replica comparisons omit setup and readback.
On the same NVIDIA L4, the current homogeneous-replica WebGPU engine delivered only 0.39--0.73 times the throughput of upstream STORMM CUDA for the recorded alanine tests.
The complete conformer workflow still lacks a matched, modern end-to-end CPU baseline.

The defensible conclusion is workload-specific:
WebGPU did not make every molecular calculation fast.
It made sufficiently large or batched local browser workloads useful enough to change the interaction design.

\subsection{Did the conformer workflow become more predictive?}

That remains unproven.
\molarium{} starts from symmetry-pruned ETKDGv3 seeds, briefly polishes them with MMFF94 or UFF, advances them together with a homogeneous-replica Sage/OBC2 workflow, and optionally refines supported structures with vacuum ANI-2x.
Final candidates can be compared with a common declared score and clustered by symmetry-aware heavy-atom RMSD.

In an early eight-molecule development panel, the replica lane found a lower common Sage minimum for four molecules.
The median molecule-level best-energy gain was only 0.128~kcal/mol.
It added 75 structural clusters, including 44 within 3~kcal/mol that were absent from the seed lane's window.
This supports increased exploration and diversity.
It does not establish generally better conformer quality.

The panel was small, used limited seeds, and judged structures with a model closely related to the replica objective.
A publication-quality test requires a frozen chemically diverse panel, repeated seeds, reference-conformer recall, symmetry-aware RMSD and torsion coverage, convergent optimization, uncertainty over molecules and seeds, matched end-to-end timing, and an independent higher-level or experimental judge for a tractable subset.
\molarium{}'s repository now states those requirements explicitly~\cite{wiewiora2026molarium}.

Table~\ref{tab:claims} summarizes the epistemic center of the case study.

\begin{table}[htbp]
\centering
\caption{Implementation, performance, and prediction are different claims.}
\label{tab:claims}
\footnotesize
\begin{tabularx}{\linewidth}{@{}>{\raggedright\arraybackslash}p{0.16\linewidth} >{\raggedright\arraybackslash}X >{\raggedright\arraybackslash}X@{}}
\toprule
\textbf{Claim} & \textbf{What the current evidence supports} & \textbf{What it does not support} \\
\midrule
\textbf{Implementation} & Selected browser kernels reproduce declared equations and reference outputs within recorded tolerances & General force-field accuracy, parameter-assignment validity, model transferability, or cross-GPU equivalence \\
\textbf{Performance} & Selected large or batched browser workloads outperform OpenMM Reference WebAssembly or earlier browser paths under stated timing boundaries & Universal speed, superiority to modern native CPU/GPU software, or an end-to-end conformer-workflow advantage \\
\textbf{Prediction} & The conformer workflow can add structural diversity and explore additional low-scoring regions & Generally better conformers, prospective drug-discovery utility, or improved molecular prediction \\
\bottomrule
\end{tabularx}
\end{table}

The ability to generate implementation rapidly is a real result.
It should not be inflated into a result about molecular truth.

\section{What does not become free}
\label{sec:not-free}

\molarium{} makes the broader argument concrete because the things it could not generate are more important than the things it could.

\paragraph{Scientific validation does not become free.}
Numerical parity requires independent energies, forces, invariances, and dynamics.
Workflow utility requires frozen benchmarks and external judges.
Predictive value requires prospective experimental evidence.
Validation is not a final ceremony applied after an agent writes code; it is the specification that makes generated scientific software meaningful.

\paragraph{Models and representations do not become free.}
Agentic orchestration can expose docking, molecular dynamics, free-energy methods, quantum chemistry, and learned potentials through a better interface.
It does not repair a force-field error, create missing biological context, solve rare-event sampling, calibrate uncertainty, or produce a more faithful representation of molecular state.
The next scientific frontier remains better models and better molecular representations, not merely better access to existing ones~\cite{bender2021part1,bender2021part2,wigh2022representation}.

\paragraph{High-quality data and proprietary knowledge do not become free.}
A discovery organization's deepest assets include assay and compound data, negative results, historical SAR, structures and biophysics, model calibration and failure modes, target-specific knowledge, medicinal-chemistry judgment, design rationale, and the reasons compounds were advanced or rejected.
This is not simply a data lake.
It is accumulated institutional intelligence.

\paragraph{Skills do not validate themselves.}
Agents make it possible to encode institutional intelligence as reusable scientific skills: prepare a protein according to an internal protocol; choose protonation states using an evidence hierarchy; analyze an RBFE network and flag suspect convergence; or balance potency, permeability, clearance, novelty, synthetic accessibility, and program objectives.
A credible skill must carry versions, provenance, domain limits, uncertainty, controls, failure conditions, tests, and an accountable owner.
Work on reusable agent-created tools and the FAIR principles for scientific objects provide part of the intellectual foundation, but scientific skills extend the requirement from reusable data or code to reusable, governed action~\cite{cai2024toolmakers,wilkinson2016fair}.

\paragraph{Security, scale, and stewardship do not become free.}
A local-first application still needs a reviewed build, dependency and model hashes, network controls, incident response, user support, and maintenance across changing browsers and hardware.
Agent-generated code can multiply faster than expert attention, amplifying the technical debt already familiar in machine-learning systems~\cite{sculley2015technicaldebt}.
Low creation cost makes ownership more---not less---important.

\paragraph{Scientific judgment does not become free.}
Automation creates the familiar irony that the human becomes most important precisely when the system leaves its normal regime~\cite{bainbridge1983ironies}.
The scarce questions are not how to click through a workflow.
They are: What result looks wrong?
What control is independent?
What approximation is acceptable?
What benchmark would change belief?
Which uncertainty should remain visible?
What should be built next?

\molarium{}'s development record is a compact demonstration.
The agent could implement a fragment builder; the chemist had to notice the flat trifluoromethyl group.
The agent could produce an energy; the chemist had to question its scale and provenance.
The agent could make a benchmark fast; the scientist had to decide whether the baseline and timing boundary were honest.
As implementation accelerates, taste, suspicion, and judgment become the bottleneck.

\section{Implications for molecular-design organizations}
\label{sec:organizations}

The commercial implication is not that scientific software disappears, nor that one public browser workbench proves an industry-wide economic transition.
It is a concrete illustration of why the center of defensible value may shift.

Generic workflow software becomes less differentiating when competitors and customers can reproduce much of it.
Validated scientific assets, proprietary evidence, secure operation, and encoded organizational judgment become more differentiating.
Vendors may increasingly expose value through trusted models, benchmarked services, data products, governed skills, and verification infrastructure, while interfaces become more dynamic and project-specific.
Internal platform teams may spend less time recreating generic screens and more time making proprietary intelligence operational.

The same shift changes the unit of software value.
The traditional unit was the application or feature.
The emerging unit may be a validated skill: a versioned decision process that combines models, tools, institutional rules, uncertainty, provenance, and acceptance tests.
Such a skill can be invoked through many interfaces and improved as evidence accumulates.
It can also be wrong, stale, or locally biased, so it requires the same governance as a model.

\molarium{} is not yet such an organizational discovery system.
It is a public case study that exposes the enabling pattern.
A scientist described the missing environment, tested it while it was being constructed, and kept the resulting implementation with its limitations and checks.
The same pattern can be applied inside a discovery organization, where the inputs are not only public equations and tools but also proprietary protocols, data, and design history.

\section{Limits of the case study}
\label{sec:limits}

\molarium{} is one retrospective, author-involved case.
It was developed by one expert scientist with one coding-agent system on a particular codebase and hardware environment.
There was no randomized comparison, matched human engineering team, blinded evaluation, or prospective time accounting.
The project was selected for discussion because it became unusually substantial.
It therefore cannot establish a general productivity effect or the fraction of commercial scientific software that can be regenerated reliably.

The current technical evidence is also developmental.
Most numerical results come from one primary browser/GPU stack; the cross-vendor matrix and frozen result ledger remain open work.
Performance comparisons use explicit but heterogeneous baselines and timing boundaries.
The conformer panel is too small to support a general prediction claim.
The provenance archive needed to distinguish generated, edited, and reused artifacts in detail is planned but not yet a finished publication package.

These limitations do not negate the case, but they determine its role.
\molarium{} should be presented as a documented case and a source of falsifiable hypotheses: that expert--agent teams can generate bounded scientific infrastructure rapidly; that validation becomes the central design problem; and that value shifts toward models, evidence, and judgment as implementation becomes easier to reproduce.
Broader claims require broader studies.

\section{Conclusion}
\label{sec:conclusion}

\molarium{} began as a request for a better place to look at and edit molecules.
It became a local browser workbench containing substantial new interaction logic, scientific workflows, WebGPU molecular mechanics, homogeneous-replica dynamics, and GPU-resident ANI-2x descriptors and forces.
The artifact is technically interesting, but its broader significance lies in what it reveals about scarcity.

An expert scientist could ask for missing software during the work itself, watch it fail, demand the relevant comparison, and retain the accepted implementation as inspectable code.
The software did not need to pre-exist as a product.
That is generated scientific infrastructure.

Nothing in the process made prediction free.
Published equations still required faithful implementation.
Parameter sets and model weights had to be preserved.
Agreement with OpenMM or TorchANI did not validate the shared science.
Performance did not establish conformer quality.
A browser workbench did not supply prospective evidence, uncertainty calibration, proprietary data, cross-platform operations, or scientific judgment.

As implementation becomes easier to reproduce, those scarce assets become more visible.
The future suggested by \molarium{} is not molecular design without scientists or without scientific-software companies.
It is molecular design in which implementation is less often the rate-limiting step---and in which better models, better data, better validation, better institutional memory, and better judgment become the clearest sources of value.

\begin{quote}
\centering
\textbf{Software implementation is becoming cheap. Scientific truth is not.}
\end{quote}

\section*{AI assistance and availability}

GPT-5.6 Sol was used as the primary coding agent during \molarium{}'s development~\cite{openai2026gpt56}.
Its outputs were directed, inspected, tested, and revised by the human project lead.
A GPT-5.6 model also assisted with restructuring and drafting this manuscript.
The authors are responsible for the scientific claims, source attribution, final text, and decision to publish.

\molarium{} source code and current technical documentation are publicly available at \url{https://github.com/rafwiewiora/molarium} and the application is available at \url{https://molarium.org}.
A frozen validation ledger and redacted development-provenance archive should accompany any submitted version that relies on the quantitative case-study claims.

\bibliographystyle{unsrtnat}
\bibliography{references}

\end{document}
